\begin{solapartat}
\punts{0,1} Primer calcularem la magnitud del camp magnètic creat pel tram recte:
\[ B_1 = \frac{\mu_0 I}{2\pi r} = \frac{4\pi\times 10^{-7}\cdot 2}{2\pi\cdot 0,02} = 2,00\times 10^{-5}\ \mathrm{T} \]
\figura[0.3]{sol_fig1.png}

\punts{0,3} Sabem que les línies de camp són circumferències en el pla $xy$ centrades en el fil, i el sentit es determina amb la mà dreta posant el dit polze en el mateix sentit que el corrent, com s'indica en la figura.
\figura[0.28]{sol_fig2.png}
Com que el camp magnètic és tangent a les línies i té el mateix sentit, el camp magnètic creat pel fil (1) té la direcció i el sentit indicats a la figura:
\[ \vec{B}_1 = -2,00\times 10^{-5}\,\vec{\imath}\ \mathrm{T} \]

\punts{0,1} Ara calcularem la magnitud del camp magnètic creat pel tram circular:
\[ B_2 = \frac{\mu_0 I}{2R} = \frac{4\pi\times 10^{-7}\cdot 2}{2\cdot 0,02} = 2\pi\times 10^{-5}\ \mathrm{T} \]
\figura[0.3]{sol_fig3.png}

\punts{0,3} La direcció del camp és perpendicular a l'espira i el sentit l'obtenim a partir de la regla de la mà dreta:
\[ \vec{B}_2 = -2\pi\times 10^{-5}\,\vec{k}\ \mathrm{T} \]

\punts{0,25} El camp magnètic total és la suma vectorial dels camps magnètics $\vec{B}_1$ i $\vec{B}_2$:
\figura[0.3]{sol_fig4.png}
\begin{align*}
\vec{B} &= \left(-2,00\times 10^{-5}\,\vec{\imath} - 2\pi\times 10^{-5}\,\vec{k}\right)\ \mathrm{T}\\
&= (-2, 0, -2\pi)\times 10^{-5}\ \mathrm{T}
\end{align*}

\punts{0,2} I finalment el mòdul és:
\[ B = \left(\sqrt{2^2 + (2\pi)^2}\right)\times 10^{-5} = 6,59\times 10^{-5}\ \mathrm{T} \]
\end{solapartat}

\begin{solapartat}
\punts{0,4} Per obtenir un camp magnètic mínim, el camp $\vec{B}_2$ ha de tenir la mateixa direcció que el camp $\vec{B}_1$ però el sentit ha de ser l'oposat:
\[ \vec{B}_2 = 2\pi\times 10^{-5}\,\vec{\imath}\ \mathrm{T} \]
\figura[0.35]{sol_fig5.png}
Per tant, aplicant la regla de la mà dreta, l'espira s'ha d'orientar en el pla $yz$ i el corrent ha de girar en sentit antihorari.

\punts{0,25} I el camp magnètic total és:
\[ \vec{B} = \left(-2,00\times 10^{-5}\,\vec{\imath} + 2\pi\times 10^{-5}\,\vec{\imath}\right)\ \mathrm{T} = 4,28\times 10^{-5}\,\vec{\imath}\ \mathrm{T} \]
I el seu mòdul és:
\[ B = 4,28\times 10^{-5}\ \mathrm{T} \]

\punts{0,4} Per obtenir un camp magnètic màxim, el camp $\vec{B}_2$ ha de tenir la mateixa direcció i sentit que el camp $\vec{B}_1$:
\[ \vec{B}_2 = -2\pi\times 10^{-5}\,\vec{\imath}\ \mathrm{T} \]
\figura[0.35]{sol_fig6.png}
Per tant, aplicant la regla de la mà dreta, l'espira s'ha d'orientar en el pla $yz$ i el corrent ha de girar en sentit horari.

\punts{0,2} I el camp magnètic total és:
\[ \vec{B} = \left(-2,00\times 10^{-5}\,\vec{\imath} - 2\pi\times 10^{-5}\,\vec{\imath}\right)\ \mathrm{T} = -8,28\times 10^{-5}\,\vec{\imath}\ \mathrm{T} \]
I el seu mòdul és:
\[ B = 8,28\times 10^{-5}\ \mathrm{T} \]
\end{solapartat}
